We used SHAP to inspect whether a ConvNeXt-Tiny classifier attends to the
correct region when it confuses visually similar ImageNet classes. Across
22 images spanning 10 confusable-class groups, the model erred on 5 (23\%),
and four of those five errors were within the intended confusable pair.
In every case we examined, SHAP attribution localized to the object itself
rather than the background, on both correct and incorrect predictions. The
distinguishing factor between the model's successes and failures appears to
be less about \emph{where} it looked and more about whether the
discriminating texture cue -- coat pattern, facial marking -- was actually
visible and salient in that region of the source image. This points to a
more specific target for improving robustness on confusable classes:
curating or augmenting training data so the discriminating cue is
consistently visible, rather than assuming the model needs to be redirected
to a different region of the image entirely.

We also resolved the specific open question our baseline resolution left
about the model's use of fine-grained cues: refining SHAP to
\texttt{max\_evals=400} on just the 5 error cases showed that our
\texttt{max\_evals=100} baseline had genuinely understated the model's
attention to a fine-grained feature in one case (the kit fox's ears), but
accurately reflected the model's attention in the other four. Prediction
errors in this study are therefore best explained by a mix of missing or
ambiguous visual evidence in the source photo and weak-confidence
tie-breaking between visually similar classes, not by a systematic failure
of the model to look at the right feature, and not merely by an artifact of
low-resolution SHAP attribution.
